\section{Six algebraic block families}\label{app:families}
For each family, $\ell$ lists the input lengths, $\pi$ lists the output
block labels, and $C$ lists the nominal sums
$C_j=1+A_j+B_j+\ell_j$ of Lemma~\ref{lem:blocks}.
Zero-length blocks are deleted from both orders; only the sums on
nonempty blocks need to be cubes. The following formulas and parameter
choices give the 25 seed lengths in Proposition~\ref{prop:finite}.

\subsection{Four, five and six blocks}
\paragraph{Family I.}
For integer cubes $U<V$, put
\[
a=2U-V-1,\qquad b=2V-3U+1.
\]
If $a,b>0$, take
\[
\ell=(a,b,a,b),\qquad \pi=(3,4,1,2),\qquad C=(U,V,U,V).
\]
The total length is $N=2(V-U)$. The three choices
$(U,V)=(64,125),(125,216),(216,343)$ give, respectively,
$(a,b)=(2,59),(33,58),(88,39)$ and $N=122,182,254$.

\paragraph{Family II.}
For integer cubes $X<Y<Z$, put
\[
a=2Y-Z-1,\qquad b=X+2Z-4Y+1,\qquad c=Y-X.
\]
If $a,b,c>0$, take
\[
\ell=(a,b,c,a,c),\qquad \pi=(4,2,5,1,3),\qquad C=(Y,X,Z,Y,Z).
\]
Here $N=2Y-X-1$. With $(Y,Z)=(64,125)$, the choices $X=27,8$
give $(a,b,c)=(2,22,37),(2,3,56)$ and $N=100,119$.

\paragraph{Family III.}
For integer cubes $X,W,Y,Z$, put
\[
\begin{split}
a&=2Y-Z-1,\qquad b=X-2a-1,\\
c&=W+2a-2X+1,\qquad d=Y+X-W-2a-1.
\end{split}
\]
If $a,b,c,d>0$, take
\[
\ell=(a,b,c,d,a,d),\qquad \pi=(5,2,3,6,1,4),\qquad C=(Y,X,W,Z,Y,Z).
\]
Then $N=2Z-2Y+X-W$. With $(Y,Z)=(64,125)$, the choices
$(X,W)=(8,64),(27,64),(8,27)$ give
$(a,b,c,d)=(2,3,53,3),(2,22,15,22),(2,3,16,40)$ and
$N=66,85,103$, respectively.

\begin{proof}
For each displayed output order, substitution in
$C_j=1+\sum_{i<j}\ell_i+
\sum_{\pp_\pi(i)<\pp_\pi(j)}\ell_i+\ell_j$
gives the indicated vector $C$ and total $N$.
All its entries are cubes and the listed lengths are positive, so
Lemma~\ref{lem:blocks} applies.
\end{proof}

\subsection{Twelve blocks with three optional blocks}
\paragraph{Family IV.}
Choose integers $X,T,Q$ and put
\[
\begin{split}
u&=X-7,\qquad v=T-2X+7,\qquad w=X+27-T,\\
z&=Q+2T-2X-121,\qquad h=X+152-T-Q.
\end{split}
\]
Assume $u,v,z\ge0$ and $w,h>0$. Set
\[
\begin{split}
\ell&=(3,u,v,w,3,w,z,h,3,27,3,h),\\
\pi&=(11,2,3,6,9,4,7,12,5,10,1,8),\\
C&=(125,X,T,64,125,64,Q,216,125,216,125,216).
\end{split}
\]
The total length is $N=276+X-T-Q$. Require $X$ to be a cube when
$u>0$, $T$ to be a cube when $v>0$, and $Q$ to be a cube when $z>0$.
Choose $X\in\{7,8\}$, then $T\in\{2X-7,27\}$, and then
$Q\in\{121-2T+2X,125\}$. The eight instances are
\[
\begin{array}{rrrrrrr}\toprule
X&T&Q&u&v&z&N\\\midrule
7&27&125&0&20&44&131\\
8&27&125&1&18&42&132\\
8&9&125&1&0&6&150\\
7&7&125&0&0&4&151\\
7&7&121&0&0&0&155\\
8&9&119&1&0&0&156\\
8&27&83&1&18&0&174\\
7&27&81&0&20&0&175\\\bottomrule
\end{array}
\]
\begin{proof}
The prefix-sum identity gives the displayed $C$ and $N$.
For these choices $w\in\{7,8,27,26\}$ and $h>0$ by substitution.
If $X=7$ then $u=0$; if $T=2X-7$ then $v=0$; and if
$Q=121-2T+2X$ then $z=0$. Otherwise the respective sums are the
cubes $8$, $27$ and $125$. Every other nominal sum is a cube.
Deleting the zero blocks therefore gives the required permutations.
\end{proof}

\subsection{Fifteen blocks with two binary parameters}
\paragraph{Family V.}
Let $\alpha,\beta\in\{0,1\}$ with $\alpha\le\beta$, put
$t_0=7+4\beta+2\alpha$, and take $t\in\{0,t_0\}$. Define
\[
\begin{aligned}
(\ell_1,\ldots,\ell_5)
 &= (2\beta,\ 7-2\alpha-4\beta,\ \alpha,\ \alpha,\ 13+4\beta),\\
(\ell_6,\ldots,\ell_{10})
 &= (\alpha,\ 23-4\beta-\alpha,\ 13+4\beta,\ 48-45\beta,\ 2\beta),\\
(\ell_{11},\ldots,\ell_{15})
 &= (\beta(39-\alpha),\ t,\ 30+\alpha-t,\ 13+4\beta,\ 91+\alpha-t),\\
\pi&=(10,4,6,2,8,3,11,1,9,14,7,12,15,5,13).
\end{aligned}
\]
Using $\alpha\beta=\alpha$, the nominal sums are
\[
\begin{aligned}
(C_1,\ldots,C_5)&=(21+43\beta,\ 8,\ 28-\alpha,\ 8,\ 216),\\
(C_6,\ldots,C_{10})&=(21+2\alpha+4\beta,\ 125,\ 64,\ 125,\ 105-41\beta),\\
(C_{11},\ldots,C_{15})&=(125,\ 209-2\alpha-4\beta+t,\ 343,\ 216,\ 343).
\end{aligned}
\]
The total length is $N=238+\alpha+2\beta-t$. Thus
\[
\begin{array}{rrrr}\toprule
\alpha&\beta&t&N\\\midrule
0&0&0&238\\0&0&7&231\\
0&1&0&240\\0&1&11&229\\
1&1&0&241\\1&1&13&228\\\bottomrule
\end{array}
\]
\begin{proof}
The lengths are nonnegative for the six listed choices.
Substitution in the prefix-sum identity gives the displayed $C$.
Blocks 1, 10 and 11 are nonempty only when $\beta=1$, with sums
$64,64,125$. Blocks 3, 4 and 6 are nonempty only when
$\alpha=\beta=1$, with sums $27,8,27$.
The sum of block 12 is needed only when $t=t_0$, and then equals $216$.
The remaining blocks have sums $8,216,125,64,125,343,216,343$.
All sums on nonempty blocks are therefore cubes, so
Lemma~\ref{lem:blocks} applies.
\end{proof}

\subsection{Ten blocks with nested binary parameters}
\paragraph{Family VI.}
Let $\alpha,\beta\in\{0,1\}$ with $\beta\le\alpha$. Set
\[
\begin{aligned}
(\ell_1,\ldots,\ell_5)
 &= (7,\ 12\alpha,\ 159-24\alpha-99\beta,\ 14\beta,\ 28\beta),\\
(\ell_6,\ldots,\ell_{10})
 &= (63\beta,\ 10+24\alpha-6\beta,\ 8\beta,\ 159-24\alpha-107\beta,
       \ 10+24\alpha+93\beta),\\
\pi&=(1,2,8,5,6,9,4,10,3,7).
\end{aligned}
\]
The nominal sums and total length are
\[
\begin{aligned}
(C_1,\ldots,C_5)&=(8,\ 15+12\alpha,\ 343,\ 333-24\alpha-93\beta,\ 174-49\beta),\\
(C_6,\ldots,C_{10})&=(174+42\beta,\ 512,\ 184+24\alpha+8\beta,\ 343,\ 512),\\
N&=345+12\alpha-6\beta.
\end{aligned}
\]
The choices $(\alpha,\beta)=(0,0),(1,0),(1,1)$ give
$N=345,357,351$, with five, six and ten nonempty blocks, respectively.
\begin{proof}
Direct substitution gives the displayed nominal sums, and all lengths are
nonnegative. Blocks 1, 3, 7, 9 and 10 are always present, with sums
$8,343,512,343,512$. Block 2 is present only at $\alpha=1$, when its
sum is $27$. Blocks 4, 5, 6 and 8 are present only at $\beta=1$,
hence $\alpha=1$, with sums $216,125,216,216$.
Lemma~\ref{lem:blocks} gives all three permutations.
\end{proof}

\section{Complete forced-neighbor rounds at the critical sizes}\label{app:peeling}
In a remaining vertex set $V$, the notation $S:c$ means that every
$i\in S$ has the unique neighbor $c-i$ in $V$. Remove $i$ and $c-i$
for all such $i$ simultaneously, counting a loop only once.
A semicolon separates groups in the same round.
All neighbor tests below use the six cubes
$\{8,27,64,125,216,343\}$; no larger cube can occur when $n\le181$.

\subsection{The first three rounds}
For $173\le n\le181$, use
\[
\begin{array}{cll}\toprule
\text{round}&S:c&\text{image interval}\\\midrule
1 &[125,342-n]:216 &[n-126,91]\\
2 &[34,215-n]:64 &[n-151,30]\\
3 &[95,276-n]:216 &[n-60,121]\\\bottomrule
\end{array}
\]
\begin{proof}
For round 1, $i\ge125$ excludes cubes at most 125, while
$343-i\ge n+1$; thus only $216-i$ is available.
For round 2, the alternative $125-i$ is in the deleted image interval
$[n-126,91]$, while $216-i>n$; cubes at most 27 are too small.
For round 3, the alternative $125-i$ lies in the deleted interval
$[n-151,30]$ and $343-i>n$; cubes at most 64 are too small.
In each round the displayed source and image intervals are disjoint
and still present. Their sizes are $218-n$, $182-n$ and $182-n$,
respectively, so these rounds remove $1164-6n$ vertices.
\end{proof}

\subsection{The remaining rounds for \texorpdfstring{$n=173$}{n=173}}
After the first three rounds, continue with the following table.
Each row is one round, not a simultaneous grouping of successive rows.
\[
\begin{array}{cl}\toprule
\text{round}&S:c\\\midrule
4 &\{4,5\}:8;\quad [8,12]:27\\
5 &\{45,46,106,107,108,122\}:216\\
6 &\{31\}:64;\quad\{172,173\}:216\\
7 &\{92\}:216\\
8 &\{1\}:8\\
9 &\{20\}:125\\
10&\{111\}:125\\
11&\{13\}:125\\
12&\{104\}:125\\
13&\{6\}:8\\
14&\{123\}:216\\
15&\{32\}:64\\\bottomrule
\end{array}
\]
At every round, subtracting each listed endpoint from the six possible
cubes shows that its only remaining neighbor is the displayed partner.
These rounds exhaust $[173]$. The loops occur at $4$, $108$ and $32$;
all other removals are reciprocal pairs. Lemma~\ref{lem:leaf} therefore
gives $P(173)=1$.

\subsection{The remaining rounds for \texorpdfstring{$n=174$ and $n=175$}{n=174 and n=175}}\label{app:175}
Round 4 is $\{4\}:8;\ [8,11]:27$ for $n=174$, and
$[8,10]:27$ for $n=175$. For both sizes the subsequent rounds are
\[
\begin{array}{cl}\toprule
\text{round}&S:c\\\midrule
5 &\{45,46,47,106,107,108\}:216\\
6 &\{15\}:27;\quad\{172,173,174\}:216\\
7 &\{113\}:216\\
8 &\{22\}:27\\
9 &\{3\}:125\\
10&\{94\}:125\\
11&\{33\}:125\\
12&\{124\}:125\\
13&\{7\}:27\\
14&\{105\}:216\\
15&\{14\}:27\\
16&\{112\}:216\\
17&\{21\}:27\\
18&\{2\}:125\\
19&\{93\}:125\\\bottomrule
\end{array}
\]
Again each listed endpoint has exactly one remaining neighbor, verified
by subtraction from the six cubes. For $n=174$ the residual set is empty,
with loops at $4$ and $108$, so $P(174)=1$.
For $n=175$, precisely the ten vertices
\[
\{4,11,16,23,41,48,102,114,168,175\}
\]
remain. Testing their pairwise sums gives the loop at $4$, the edge
$4\sim23$, and the chordless cycle
\[
23\sim41\sim175\sim168\sim48\sim16\sim11\sim114\sim102\sim23.
\]
Its consecutive edge sums are $64,216,343,216,64,27,125,216,125$.
If $4$ is fixed, the odd cycle has the two orientations as its only
permutation covers: reciprocal pairs cannot cover an odd number of vertices.
If $4$ is paired with $23$, the remaining path has a unique cover by
four reciprocal pairs. Thus the residual graph has exactly three
permutations, and Lemma~\ref{lem:leaf} gives $P(175)=3$.
